% ============================================================
%  cap05_saude.tex — Capítulo 5: Saúde
% ============================================================

% ════════════════════════════════════════════════════════════
\chapter{Saúde}
% ════════════════════════════════════════════════════════════

\begin{boxdestaque}{Neste capítulo}
Este capítulo apresenta as principais fontes de dados sobre saúde disponíveis para o Brasil. Cobrimos o DATASUS e seus sistemas componentes --- SIM, SINASC, SIH, SIA, SINAN, CNES, SIOPS, SISVAN e SISAB ---, a Pesquisa Nacional de Saúde (PNS) e a Pesquisa de Assistência Médico-Sanitária (AMS). Para cada fonte, detalhamos histórico, estrutura, variáveis disponíveis, forma de acesso e aplicações típicas em avaliação de políticas de saúde pública.
\end{boxdestaque}

% ────────────────────────────────────────────────────────────
\section{Introdução: o ecossistema de dados de saúde no Brasil}
% ────────────────────────────────────────────────────────────

O Brasil dispõe de um dos sistemas de informação em saúde mais abrangentes do mundo em desenvolvimento, construído ao longo de décadas a partir da estrutura do Sistema Único de Saúde (SUS). O DATASUS --- Departamento de Informática do SUS, criado em 1991 --- é o principal repositório desses dados, reunindo em uma única plataforma sistemas de mortalidade, nascimentos, internações, ambulatório, vigilância epidemiológica, cadastro de estabelecimentos e orçamento público em saúde.

Essa riqueza de fontes é um ativo extraordinário para a avaliação de políticas de saúde. Ao mesmo tempo, o ecossistema é complexo: cada sistema tem sua própria lógica de produção, periodicidade, cobertura e qualidade. Entender o que cada sistema mede --- e, igualmente importante, o que não mede --- é o ponto de partida para qualquer análise séria.

A Figura~\ref{fig:datasus_sistemas} apresenta uma visão esquemática dos principais sistemas do DATASUS e sua relação com os diferentes pontos de contato do cidadão com o sistema de saúde.

\begin{table}[H]
\centering
\caption{Principais sistemas do DATASUS e suas coberturas}
\label{fig:datasus_sistemas}
\small
\begin{tabularx}{\textwidth}{@{} p{1.8cm} p{3.5cm} >{\raggedright\arraybackslash}p{2.5cm} X @{}}
\toprule
\textbf{Sistema} & \textbf{O que registra} & \textbf{Desde} & \textbf{Principal uso em avaliação} \\
\midrule
SIM    & Óbitos                        & 1979 & Mortalidade, causas de morte, violência \\
SINASC & Nascimentos vivos             & 1990 & Saúde materna, prematuridade, desigualdades \\
SIH    & Internações hospitalares SUS  & 1992 & Uso hospitalar, procedimentos, custos \\
SIA    & Atendimentos ambulatoriais SUS & 1994 & Produção ambulatorial, cobertura \\
SINAN  & Agravos de notificação        & 1993 & Doenças infecciosas, vigilância epidemiológica \\
CNES   & Estabelecimentos de saúde     & 2005 & Infraestrutura, RH, capacidade instalada \\
SIOPS  & Orçamento e execução em saúde    & 2000 & Gasto público, financiamento do SUS \\
SISVAN & Vigilância alimentar/nutricional & 2008 & Desnutrição, obesidade, segurança alimentar \\
SISAB  & Atenção primária à saúde      & 2013 & Cobertura de AB, ESF, ações de saúde na AB \\
\bottomrule
\end{tabularx}
\fonte{FGV CLEAR, elaboração própria com base em DATASUS.}
\end{table}

% ────────────────────────────────────────────────────────────
\section{SIM --- Sistema de Informações sobre Mortalidade}
% ────────────────────────────────────────────────────────────

\ficharapida
  {SIM --- Sistema de Informações sobre Mortalidade}
  {Ministério da Saúde / DATASUS}
  {https://datasus.saude.gov.br/mortalidade-desde-1996-pela-cid-10}
  {Brasil, Unidades da Federação e municípios}
  {Anual (desde 1979; dados consistentes a partir de 1996 com a CID-10)}
  {Óbito}

\subsection{O que é e por que importa}

O Sistema de Informações sobre Mortalidade (SIM) é o registro nacional de óbitos no Brasil. Criado em 1979, o SIM coleta informações de todas as Declarações de Óbito (DO) emitidas no país, processadas pelas Secretarias Estaduais de Saúde e consolidadas pelo Ministério da Saúde. A partir de 1996, com a adoção da Classificação Internacional de Doenças em sua décima revisão (CID-10), as séries do SIM tornaram-se mais comparáveis internacionalmente e consistentes para análises de longo prazo.

Para a avaliação de políticas públicas, o SIM é uma fonte insubstituível porque a morte é um desfecho objetivo, verificável e registrado de forma relativamente padronizada --- ao contrário de muitos outros indicadores de saúde que dependem de autoavaliação ou de acesso a serviços. Políticas de saúde materno-infantil, de atenção primária, de combate à violência e de doenças crônicas têm no SIM uma das principais fontes de avaliação de resultados e impacto.

\subsection{Estrutura e principais variáveis}

Cada registro do SIM corresponde a uma Declaração de Óbito e contém:

\begin{itemize}
  \item \textbf{Características do falecido}: sexo, idade, raça/cor, escolaridade, estado civil e ocupação --- a qualidade de preenchimento dessas variáveis é heterogênea; em particular, escolaridade, raça/cor e ocupação apresentam taxas expressivas de \textit{missing} e exigem tratamento cuidadoso antes da análise;
  \item \textbf{Causa da morte}: causa básica e causas associadas, codificadas pela CID-10 (adotada no Brasil em 1996, com implementação plena nos registros do SIM a partir de 1998; séries anteriores usam a CID-9 e exigem compatibilização para análises históricas); causa externa (para mortes violentas);
  \item \textbf{Localização}: município de residência e município de ocorrência do óbito;
  \item \textbf{Data e local do óbito}: data, hora, local (hospital, domicílio, via pública, etc.);
  \item \textbf{Para óbitos infantis e fetais}: informações adicionais sobre a mãe (idade, número de filhos, gestação) e sobre o recém-nascido (peso ao nascer, duração da gestação).
\end{itemize}

\subsection{Aplicações típicas em avaliação}

\begin{itemize}
  \item \textbf{Mortalidade infantil e materna}: o SIM, combinado com o SINASC (para os denominadores de nascidos vivos), é a fonte padrão para calcular as taxas de mortalidade infantil e materna por município;
  \item \textbf{Atlas da Violência}: o Instituto de Pesquisa Econômica Aplicada (Ipea) e o Fórum Brasileiro de Segurança Pública produzem anualmente o Atlas da Violência a partir de mortes por causas externas (homicídios, acidentes de trânsito) registradas no SIM;
  \item \textbf{Avaliação do Programa Saúde da Família}: um dos usos mais citados do SIM em avaliação de impacto é a estimação dos efeitos da expansão do PSF sobre mortalidade infantil e materna em municípios brasileiros;
  \item \textbf{Doenças crônicas não transmissíveis}: mortalidade por doenças cardiovasculares, câncer e diabetes para avaliação de programas de rastreamento e tratamento.
\end{itemize}

% CORRIGIDO: título encurtado para caber na caixa
\begin{boxatencao}{Atenção --- Subregistro e causa básica no SIM}
O SIM tem dois problemas de qualidade conhecidos que o analista deve considerar. Primeiro, o \textbf{subregistro de óbitos}: historicamente, uma parcela dos óbitos --- especialmente em municípios pequenos, áreas rurais e regiões Norte e Nordeste --- não era registrada, subestimando a mortalidade real nessas áreas. Esse quadro melhorou substancialmente nas últimas duas décadas, com a expansão da Estratégia Saúde da Família, iniciativas como a Busca Ativa de Óbitos e a qualificação das equipes municipais de vigilância; ainda assim, desigualdades regionais persistem em municípios menores e mais isolados. Segundo, a \textbf{qualidade do preenchimento da causa básica}: uma proporção dos óbitos é classificada como ``causas mal definidas'' (capítulo XVIII da CID-10), especialmente em regiões com menor cobertura de serviços de saúde --- proporção que também vem caindo de forma expressiva ao longo do tempo. Análises de mortalidade por causa específica devem sempre verificar a proporção de causas mal definidas no universo analisado e, se necessário, aplicar redistribuição proporcional.
\end{boxatencao}

\subsection{Acesso aos dados}

Os microdados do SIM são disponibilizados gratuitamente pelo DATASUS em formato DBC (comprimido) ou DBF. O portal TabNet (\url{tabnet.datasus.gov.br}) permite tabulações online sem necessidade de download. Para análises com microdados, a plataforma \textit{Base dos Dados} disponibiliza o SIM já tratado e compatibilizado em formato Parquet/BigQuery.

\begin{boxdica}{Dica --- Pacote \texttt{microdatasus} no R}
O pacote \texttt{microdatasus} (disponível no CRAN) facilita enormemente o download e a leitura dos microdados do SIM e de outros sistemas do DATASUS diretamente no R, convertendo automaticamente os arquivos DBC para formato utilizável e aplicando os dicionários de variáveis.
\end{boxdica}

% ────────────────────────────────────────────────────────────
\section{SINASC --- Sistema de Informações sobre Nascidos Vivos}
% ────────────────────────────────────────────────────────────

\ficharapida
  {SINASC --- Sistema de Informações sobre Nascidos Vivos}
  {Ministério da Saúde / DATASUS}
  {https://datasus.saude.gov.br/nascidos-vivos-desde-1994}
  {Brasil, Unidades da Federação e municípios}
  {Anual (desde 1990; dados mais consistentes a partir de 1994)}
  {Nascimento vivo}

\subsection{O que é e principais variáveis}

O SINASC registra todos os nascimentos vivos ocorridos no Brasil a partir das Declarações de Nascido Vivo (DNV) preenchidas nas maternidades e unidades de saúde. É a principal fonte de dados sobre saúde perinatal e materna no país.

Cada registro contém:

\begin{itemize}
  \item \textbf{Características do recém-nascido}: sexo, peso ao nascer, duração da gestação (semanas), índice de Apgar, presença de anomalia congênita;
  \item \textbf{Características da mãe}: idade, escolaridade, raça/cor, município de residência, número de filhos anteriores (nascidos vivos e mortos), número de consultas de pré-natal, tipo de parto (vaginal ou cesáreo);
  \item \textbf{Local do nascimento}: município de ocorrência, estabelecimento de saúde (quando aplicável) e local do parto (hospitalar, domiciliar, outro estabelecimento de saúde, via pública, outros).
\end{itemize}

\subsection{Aplicações típicas em avaliação}

\begin{itemize}
  \item \textbf{Denominador para taxas de mortalidade infantil}: o SINASC fornece o número de nascidos vivos por município, necessário para calcular a taxa de mortalidade infantil com dados do SIM;
  \item \textbf{Avaliação de políticas de pré-natal}: o número de consultas de pré-natal registrado no SINASC permite avaliar a cobertura e a qualidade do pré-natal por município;
  \item \textbf{Desigualdades no nascimento}: o cruzamento de variáveis como peso ao nascer, prematuridade e características socioeconômicas da mãe permite análises de equidade em saúde materna;
  \item \textbf{Impacto de políticas sobre cesáreas}: o SINASC é a fonte de referência para monitorar a proporção de partos cesáreos --- indicador de uso excessivo de intervenção cirúrgica no parto.
\end{itemize}

% ────────────────────────────────────────────────────────────
\section{SIH e SIA --- Internações e Atendimento Ambulatorial}
% ────────────────────────────────────────────────────────────

\ficharapida
  {SIH/SIA --- Sistemas de Informações Hospitalares e Ambulatoriais}
  {Ministério da Saúde / DATASUS}
  {https://datasus.saude.gov.br/producao-hospitalar-desde-1994}
  {Brasil, Unidades da Federação e municípios}
  {Mensal (SIH desde 1992; SIA desde 1994)}
  {Autorização de Internação Hospitalar (SIH) e procedimento ambulatorial (SIA)}

\subsection{O que são e o que cobrem}

O Sistema de Informações Hospitalares (SIH) e o Sistema de Informações Ambulatoriais (SIA) registram a \textbf{produção de serviços de saúde} financiados pelo SUS --- internações hospitalares e atendimentos ambulatoriais, respectivamente. Não cobrem serviços prestados pela saúde suplementar (planos de saúde privados) nem atendimentos particulares.

O SIH registra cada internação por meio da Autorização de Internação Hospitalar (AIH), que contém:
\begin{itemize}
  \item Diagnóstico principal e secundário (CID-10);
  \item Procedimento realizado;
  \item Tempo de permanência;
  \item Valor pago pelo SUS;
  \item Características do paciente (sexo, idade, município de residência);
  \item Estabelecimento de saúde (CNES).
\end{itemize}

O SIA registra cada procedimento ambulatorial realizado em estabelecimentos credenciados ao SUS, com informações sobre o tipo de procedimento, especialidade, estabelecimento e município.

\subsection{Aplicações típicas e limitações}

Os sistemas SIH e SIA são especialmente úteis para:
\begin{itemize}
  \item Avaliar \textbf{mudanças no padrão de utilização hospitalar} decorrentes de políticas de atenção primária (como a expansão do PSF, que teoricamente deveria reduzir internações evitáveis);
  \item Calcular \textbf{custos de procedimentos} para análises de custo-efetividade;
  \item Monitorar a \textbf{produção de serviços} de saúde por região e estabelecimento.
\end{itemize}

A limitação central é que SIH e SIA cobrem \textbf{apenas o SUS}: aproximadamente 25\% da população brasileira possui plano de saúde privado e seus atendimentos não constam nessas bases. Para análises da população como um todo, é preciso combinar dados do SUS com fontes que cubram a saúde suplementar.

% ────────────────────────────────────────────────────────────
\section{SINAN --- Sistema de Informação de Agravos de Notificação}
% ────────────────────────────────────────────────────────────

\ficharapida
  {SINAN --- Sistema de Informação de Agravos de Notificação}
  {Ministério da Saúde / DATASUS}
  {https://datasus.saude.gov.br/doencas-e-agravos-de-notificacao-de-2007-em-diante-sinan}
  {Brasil, Unidades da Federação e municípios}
  {Contínuo (desde 1993; integração com VIVA a partir de 2009)}
  {Caso notificado de agravo ou doença de notificação compulsória}

\subsection{O que é e principais doenças cobertas}

O SINAN reúne as notificações compulsórias de doenças e agravos à saúde --- ou seja, condições que, por lei, os profissionais e serviços de saúde são obrigados a comunicar às autoridades sanitárias. A lista de doenças de notificação compulsória é definida pelo Ministério da Saúde e inclui atualmente mais de 70 agravos.

Entre os mais relevantes para avaliação de políticas:

\begin{itemize}
  \item \textbf{Doenças infecciosas e parasitárias}: dengue, chikungunya, zika, malária, tuberculose, hanseníase, leishmaniose, hepatites virais;
  \item \textbf{Doenças imunopreveníveis}: sarampo, rubéola, coqueluche, poliomielite;
  \item \textbf{HIV/AIDS}: casos novos, perfil epidemiológico;
  \item \textbf{Violências e acidentes} (via integração com o VIVA): violência doméstica, sexual, autoprovocada, acidentes de trabalho.
\end{itemize}

\subsection{Aplicações típicas em avaliação}

\begin{itemize}
  \item \textbf{Avaliação de programas de vacinação}: a cobertura vacinal pode ser combinada com dados do SINAN sobre incidência das doenças-alvo para avaliar efetividade;
  \item \textbf{Políticas de controle de endemias}: monitoramento de dengue, malária e leishmaniose em resposta a intervenções de saneamento e controle vetorial;
  \item \textbf{Políticas de combate à tuberculose}: o SINAN é a referência para avaliação de programas de tratamento e controle da tuberculose, incluindo taxas de cura, abandono e resistência;
  \item \textbf{Violência doméstica e sexual}: desde a integração com o VIVA, o SINAN passou a ser uma das principais fontes sobre violência atendida em serviços de saúde.
\end{itemize}

% CORRIGIDO: título encurtado para caber na caixa
\begin{boxatencao}{Atenção --- Subnotificação no SINAN}
A notificação compulsória é obrigatória por lei, mas o cumprimento da obrigação varia enormemente entre doenças, regiões e tipos de serviço. Estima-se que para dengue, por exemplo, apenas uma fração dos casos reais é notificada --- especialmente em períodos de alta demanda, quando os serviços ficam sobrecarregados. Para doenças com maior estigma social (como HIV/AIDS e hepatites), a subnotificação pode ser ainda maior. Análises com dados do SINAN devem sempre discutir o potencial de subnotificação e, quando possível, triangular com outras fontes.
\end{boxatencao}

% ────────────────────────────────────────────────────────────
\section{CNES --- Cadastro Nacional de Estabelecimentos de Saúde}
% ────────────────────────────────────────────────────────────

\ficharapida
  {CNES --- Cadastro Nacional de Estabelecimentos de Saúde}
  {Ministério da Saúde / DATASUS}
  {https://cnes.datasus.gov.br}
  {Brasil, Unidades da Federação e municípios}
  {Contínuo (atualização mensal; desde 2005)}
  {Estabelecimento de saúde}

\subsection{O que é e principais variáveis}

O CNES é o cadastro de todos os estabelecimentos de saúde do Brasil --- públicos e privados, com e sem fins lucrativos, vinculados ou não ao SUS. Cada estabelecimento é identificado por um código único (CNES) que serve de chave de integração com os sistemas SIH, SIA e SINAN.

As principais variáveis disponíveis incluem:

\begin{itemize}
  \item \textbf{Tipo e natureza do estabelecimento}: hospital, UBS, UPA, clínica, laboratório, etc.; público, privado, filantrópico;
  \item \textbf{Localização}: município, UF, endereço georreferenciado;
  \item \textbf{Infraestrutura}: número de leitos por tipo (SUS, não-SUS, UTI), equipamentos disponíveis;
  \item \textbf{Recursos humanos}: profissionais de saúde vinculados por categoria (médicos, enfermeiros, técnicos), carga horária e vínculo;
  \item \textbf{Serviços oferecidos}: especialidades, procedimentos habilitados, serviços de apoio diagnóstico.
\end{itemize}

\subsection{Aplicações típicas em avaliação}

\begin{itemize}
  \item \textbf{Avaliação de acesso a serviços de saúde}: o CNES permite calcular indicadores de disponibilidade de serviços (leitos por mil habitantes, médicos por habitante) por município, essenciais para análises de equidade no acesso;
  \item \textbf{Avaliação de políticas de expansão da rede}: monitoramento da abertura ou fechamento de estabelecimentos em resposta a políticas de investimento em infraestrutura;
  \item \textbf{Variável de controle em avaliações de impacto}: características da oferta de saúde (número de leitos, cobertura de ESF) são frequentemente usadas como variáveis de controle em estudos que avaliam outros desfechos de saúde.
\end{itemize}

% ────────────────────────────────────────────────────────────
\section{SIOPS --- Sistema de Informações sobre Orçamentos Públicos em Saúde}
% ────────────────────────────────────────────────────────────

\ficharapida
  {SIOPS}
  {Ministério da Saúde}
  {https://siops.datasus.gov.br}
  {Brasil, Unidades da Federação e municípios}
  {Anual (desde 2000)}
  {Ente federativo (município, estado ou União)}

\subsection{O que é e aplicações}

O SIOPS reúne dados sobre o orçamento público destinado à saúde pelos municípios, estados e União, conforme a obrigação constitucional de aplicação mínima de recursos próprios na área. Cada ente declarante informa suas receitas totais e os gastos realizados em ações e serviços públicos de saúde (ASPS).

Para avaliação de políticas, o SIOPS é fundamental para:
\begin{itemize}
  \item Analisar o \textbf{nível de financiamento público da saúde} por município e UF ao longo do tempo;
  \item Calcular \textbf{gasto per capita em saúde} e comparar entre entes federativos;
  \item Avaliar o cumprimento dos \textbf{pisos constitucionais} de gasto em saúde (municípios devem aplicar mínimo de 15\% da receita própria; estados, 12\%).
\end{itemize}

% ────────────────────────────────────────────────────────────
\section{SISVAN --- Sistema de Vigilância Alimentar e Nutricional}
% ────────────────────────────────────────────────────────────

\ficharapida
  {SISVAN}
  {Ministério da Saúde / DATASUS}
  {https://sisvan.saude.gov.br}
  {Brasil, Unidades da Federação e municípios}
  {Contínuo (desde 2008 com cobertura ampliada)}
  {Indivíduo atendido na atenção básica}

\subsection{O que é e aplicações}

O SISVAN registra o estado nutricional e o consumo alimentar de indivíduos acompanhados pela atenção básica do SUS. Os dados são coletados nas consultas de rotina nas Unidades Básicas de Saúde e cobrem principalmente crianças menores de 5 anos, gestantes, mulheres em idade fértil e idosos.

As principais aplicações em avaliação incluem:
\begin{itemize}
  \item Monitoramento de \textbf{desnutrição e obesidade} por município, especialmente em populações vulneráveis atendidas pelo SUS;
  \item Avaliação de \textbf{políticas de segurança alimentar} e programas de suplementação nutricional;
  \item Acompanhamento nutricional de \textbf{beneficiários do Bolsa Família}, que têm condicionalidades de saúde vinculadas ao SISVAN.
\end{itemize}

A principal limitação do SISVAN é que cobre apenas indivíduos que comparecem às consultas na atenção básica --- portanto, não é representativo da população geral.

% ────────────────────────────────────────────────────────────
\section{SISAB --- Sistema de Informação em Saúde para a Atenção Básica}
% ────────────────────────────────────────────────────────────

\ficharapida
  {SISAB}
  {Ministério da Saúde / SAPS}
  {https://sisab.saude.gov.br}
  {Brasil, Unidades da Federação, municípios e equipes}
  {Contínuo (desde 2013, substituindo o SIAB)}
  {Indivíduo atendido, equipe e estabelecimento de atenção básica}

\subsection{O que é e aplicações}

O SISAB consolida a produção das equipes de atenção básica do SUS --- Estratégia Saúde da Família (ESF), Núcleos Ampliados de Saúde da Família (eMulti, sucessor do NASF-AB) e demais modelos de Unidade Básica de Saúde (UBS). Substituiu o antigo SIAB em 2013 e é alimentado pelo prontuário eletrônico do cidadão (PEC) e por formulários de CDS (Coleta de Dados Simplificada).

Para a avaliação de políticas públicas, o SISAB permite:
\begin{itemize}
  \item Acompanhar a \textbf{cobertura populacional da atenção básica} por município e por equipe, variável-chave em avaliações de impacto da ESF sobre desfechos de saúde;
  \item Construir indicadores de \textbf{produção de ações da AB} (consultas, visitas domiciliares, procedimentos, ações coletivas, visitas de ACS);
  \item Cruzar dados da atenção básica a sistemas hospitalares (SIH) e de mortalidade (SIM) via código CNES, permitindo análises do efeito da AB sobre hospitalizações sensíveis à atenção primária e mortalidade evitável.
\end{itemize}

Como no SISVAN, a limitação central é a cobertura restrita aos usuários efetivos da atenção básica e a heterogeneidade municipal no preenchimento.

% ────────────────────────────────────────────────────────────
\section{Pesquisa Nacional de Saúde (PNS)}
% ────────────────────────────────────────────────────────────

\ficharapida
  {Pesquisa Nacional de Saúde (PNS)}
  {IBGE em parceria com o Ministério da Saúde / Fiocruz}
  {https://www.ibge.gov.br/estatisticas/sociais/saude/9160-pesquisa-nacional-de-saude.html}
  {Brasil, Grandes Regiões, Unidades da Federação e capitais}
  {Quinquenal: 2013 e 2019}
  {Domicílio e morador}

\subsection{O que é e por que importa}

A Pesquisa Nacional de Saúde (PNS) é o principal inquérito domiciliar de saúde do Brasil, realizado pelo IBGE em parceria com o Ministério da Saúde e a Fiocruz. Diferentemente dos sistemas do DATASUS --- que captam apenas quem utiliza serviços de saúde ---, a PNS é uma pesquisa amostral representativa de toda a população residente em domicílios particulares, o que permite estimar a prevalência de condições de saúde na população geral, independentemente do uso de serviços.

\subsection{Principais temas cobertos}

A PNS abrange um conjunto amplo de dimensões da saúde individual e coletiva:

\begin{itemize}
  \item \textbf{Percepção do estado de saúde}: autoavaliação da saúde, limitações funcionais;
  \item \textbf{Doenças crônicas não transmissíveis}: hipertensão, diabetes, doenças cardíacas, AVC, câncer, problemas respiratórios crônicos;
  \item \textbf{Saúde mental}: depressão, ansiedade (edição 2019 com módulo específico);
  \item \textbf{Acesso e uso de serviços de saúde}: consultas médicas, internações, uso de plano de saúde, acesso à atenção básica;
  \item \textbf{Fatores de risco e proteção}: tabagismo, consumo de álcool, atividade física, alimentação;
  \item \textbf{Saúde da mulher}: mamografia, preventivo, pré-natal;
  \item \textbf{Características do domicílio e socioeconômicas}: escolaridade, renda, trabalho.
\end{itemize}

\subsection{Aplicações típicas em avaliação}

\begin{itemize}
  \item \textbf{Estimação de prevalências populacionais} de doenças crônicas e fatores de risco, impossível com os sistemas do DATASUS;
  \item \textbf{Avaliação de acesso a serviços}: comparação entre grupos populacionais (por renda, raça, região) no uso de consultas médicas, exames e internações;
  \item \textbf{Linha de base para políticas de saúde}: as edições de 2013 e 2019 fornecem um retrato do estado de saúde da população antes e depois de mudanças de política relevantes;
  \item \textbf{Desigualdades em saúde}: a PNS permite analisar gradientes socioeconômicos em saúde de forma muito mais rica do que as fontes administrativas.
\end{itemize}

\begin{boxatencao}{Atenção --- Periodicidade quinquenal}
A PNS é realizada a cada cinco anos, o que limita seu uso para avaliações de políticas de curto prazo. Entre edições, os sistemas do DATASUS são as principais fontes para acompanhamento de indicadores de saúde. A edição 2024/25 estava em planejamento no momento da elaboração deste guia; recomenda-se verificar o site do IBGE para informações atualizadas.
\end{boxatencao}

% ────────────────────────────────────────────────────────────
\section{AMS --- Pesquisa de Assistência Médico-Sanitária}
% ────────────────────────────────────────────────────────────

\ficharapida
  {AMS --- Pesquisa de Assistência Médico-Sanitária}
  {Instituto Brasileiro de Geografia e Estatística (IBGE)}
  {https://www.ibge.gov.br/estatisticas/sociais/saude/9067-pesquisa-de-assistencia-medico-sanitaria.html}
  {Brasil, Unidades da Federação e municípios}
  {Cross-sections: 1999, 2002, 2005, 2009 (descontinuada; substituída parcialmente pelo CNES)}
  {Estabelecimento de saúde}

\subsection{O que é e aplicações}

A AMS era um censo de todos os estabelecimentos de saúde do Brasil --- públicos e privados ---, realizado pelo IBGE. Diferentemente do CNES (que é um cadastro administrativo atualizado mensalmente), a AMS era uma pesquisa estruturada com questionários padronizados aplicados a cada estabelecimento, coletando informações detalhadas sobre infraestrutura, equipamentos, recursos humanos e serviços oferecidos.

Com a consolidação do CNES como fonte administrativa de referência, a AMS foi descontinuada após 2009. Para análises históricas da oferta de serviços de saúde no período 1999--2009, a AMS ainda é uma fonte relevante, especialmente para comparações que incluam estabelecimentos privados não vinculados ao SUS --- categoria que o CNES cobre com menor profundidade analítica.

% ────────────────────────────────────────────────────────────
\section{Síntese do capítulo}
% ────────────────────────────────────────────────────────────

Este capítulo percorreu o ecossistema de dados de saúde disponível para o Brasil. A Tabela~\ref{tab:sintese_cap5} resume as características essenciais de cada fonte.

\begin{table}[H]
\centering
\caption{Síntese das fontes de dados de saúde}
\label{tab:sintese_cap5}
\small
\begin{tabularx}{\textwidth}{@{} p{2cm} p{3.5cm} >{\raggedright\arraybackslash}p{2.5cm} X @{}}
\toprule
\textbf{Fonte} & \textbf{O que registra} & \textbf{Cobertura} & \textbf{Principal uso em avaliação} \\
\midrule
SIM    & Óbitos                     & Total (com subregistro) & Mortalidade, causas externas, violência \\
SINASC & Nascimentos vivos          & Total (com subregistro) & Saúde materna, prematuridade \\
SIH    & Internações SUS            & SUS apenas & Uso hospitalar, procedimentos \\
SIA    & Atendimentos ambulatoriais SUS & SUS apenas & Produção ambulatorial \\
SINAN  & Agravos notificáveis       & Notificados (subnotificação) & Vigilância epidemiológica \\
CNES   & Estabelecimentos de saúde  & Todos (públicos e privados) & Infraestrutura, acesso \\
SIOPS  & Orçamento e execução em saúde & Entes federativos & Financiamento do SUS \\
SISVAN & Estado nutricional         & Usuários da atenção básica & Nutrição, segurança alimentar \\
SISAB  & Produção da atenção básica & Usuários da atenção básica & Cobertura da ESF, ações da AB \\
PNS    & Saúde da população geral   & Amostra representativa & Prevalências, acesso, desigualdades \\
AMS    & Estabelecimentos (histórico) & Todos (descontinuada) & Séries históricas 1999--2009 \\
\bottomrule
\end{tabularx}
\fonte{FGV CLEAR, elaboração própria.}
\end{table}

Os pontos centrais do capítulo são:

\begin{itemize}
  \item O \textbf{SIM} e o \textbf{SINASC} são as fontes de referência para mortalidade e nascimentos, com cobertura nacional e séries históricas longas, mas sujeitos a subregistro diferencial por região;
  \item O \textbf{SIH} e o \textbf{SIA} registram a produção de serviços do SUS, mas não cobrem a saúde suplementar privada;
  \item O \textbf{SINAN} é fundamental para vigilância epidemiológica, mas a subnotificação é estrutural e varia muito entre agravos;
  \item O \textbf{CNES} é a referência para infraestrutura e recursos humanos em saúde, com atualização mensal;
  \item A \textbf{PNS} é a única fonte que permite estimar prevalências de condições de saúde na população geral, independentemente do uso de serviços --- mas sua periodicidade quinquenal limita o acompanhamento de curto prazo;
  \item A \textbf{combinação estratégica} das fontes --- por exemplo, SIM + SINASC para mortalidade infantil, ou CNES + SIH para análises de acesso e uso --- é o que permite avaliações mais robustas do setor saúde.
\end{itemize}

\medskip

O próximo capítulo apresenta as fontes sobre educação: o Censo Escolar, o SAEB, o ENEM, o Censo da Educação Superior e o SIOPE --- o ecossistema coordenado pelo Inep, que documenta matrículas, desempenho cognitivo, acesso ao ensino e financiamento da educação básica.